%% appA --- The C# to Python cheat sheet
%%
%% Written September 2026. One table per part of the book. Every row names
%% the chapter that TEACHES it, with \ref rather than a typed number, and
%% tools/check_structure.py --cheatsheet fails when a row points at a label
%% that does not exist or at a chapter that is still a stub. A row nothing
%% teaches is a defect in this appendix rather than a fact about the book.

\chapter{The C\# to Python cheat sheet}
\label{app:cheatsheet}

\noindent
Everything here is taught somewhere in the book, and the last column says
where. That is the point of the column: a row is a claim about a chapter,
and a row no chapter teaches is a defect in this appendix rather than a
gap in your knowledge.

Read the two middle columns as a translation and not as an equivalence.
Several rows map cleanly; several map with a difference that costs
something, and those are the ones the chapter spends its pages on. Where a
row would mislead on its own, it says so in the Python column.

\section{Part I \dash{} Runtime and toolchain}
\label{sec:appA-part1}

\begin{cheatsheet}{C\#}{Python}{Ch.}
the CLR, JIT to native & CPython, bytecode in \code{\_\_pycache\_\_}; no
  JIT to rely on & \ref{ch:runtime} \\
IL, \code{ildasm} & bytecode, \api{dis.dis} on a function &
  \ref{ch:runtime} \\
threads scale CPU work & they do not: the lock serialises bytecode. They
  scale I/O work & \ref{ch:runtime} \\
\code{static void Main} & \code{if \_\_name\_\_ == "\_\_main\_\_"} &
  \ref{ch:runtime} \\
\code{Parallel.For} & \api{multiprocessing}, or the free-threaded build &
  \ref{ch:runtime} \\
NuGet & PyPI, and \pkg{uv} to talk to it & \ref{ch:packaging} \\
\code{packages.lock.json} & \code{uv.lock}, committed &
  \ref{ch:packaging} \\
\code{dotnet restore} & \code{uv sync --locked} & \ref{ch:packaging} \\
\code{.csproj} & \code{pyproject.toml} & \ref{ch:packaging} \\
\code{dotnet run} & \code{uv run}, which resolves the environment per
  invocation & \ref{ch:packaging} \\
a global tool & a per-project environment; there is no good global install
  & \ref{ch:packaging} \\
the compiler checks types & a checker does, if you invite one: pyright or
  mypy & \ref{ch:typing} \\
\code{int?}, \code{string?} & \code{int | None}, \code{str | None} &
  \ref{ch:typing} \\
\code{List<int>}, \code{Dictionary<K,V>} & \code{list[int]},
  \code{dict[K, V]} & \ref{ch:typing} \\
an interface & \code{Protocol}, satisfied structurally, with no
  declaration & \ref{ch:typing} \\
\code{where T : IComparable} & a bound on a type parameter, PEP~695
  syntax & \ref{ch:typing} \\
\code{dynamic} & \code{Any} \dash{} but it propagates, and silences the
  checker downstream & \ref{ch:typing} \\
\code{is T t} pattern & \api{isinstance}, and \code{TypeIs} to tell the
  checker & \ref{ch:typing} \\
\end{cheatsheet}

\section{Part II \dash{} The language, mapped}
\label{sec:appA-part2}

\begin{cheatsheet}{C\#}{Python}{Ch.}
\code{Equals} + \code{GetHashCode} & \code{\_\_eq\_\_} +
  \code{\_\_hash\_\_}, and omitting the second removes hashing &
  \ref{ch:objects} \\
\code{ToString()} & \code{\_\_repr\_\_}, and optionally
  \code{\_\_str\_\_} & \ref{ch:objects} \\
operator overloading & \code{\_\_add\_\_}, \code{\_\_lt\_\_} and the rest,
  none declared & \ref{ch:objects} \\
\code{IComparable} & \code{\_\_lt\_\_}, or \code{order=True} on a
  dataclass & \ref{ch:objects} \\
\code{record}, \code{with} & \code{@dataclass(frozen=True)},
  \api{dataclasses.replace} & \ref{ch:objects} \\
a property & \code{@property}, promoted from a plain attribute when you
  need it & \ref{ch:objects} \\
an indexer, \code{Count} & \code{\_\_getitem\_\_}, \code{\_\_len\_\_} &
  \ref{ch:objects} \\
\code{IEnumerable<T>} & \code{\_\_iter\_\_} & \ref{ch:objects} \\
\code{enum}, \code{[Flags]} & \code{Enum}, \code{StrEnum}, \code{Flag} &
  \ref{ch:objects} \\
a \code{static} field & a class attribute \dash{} but readable through
  \code{self}, which is the trap & \ref{ch:objects} \\
\code{struct} & \code{\_\_slots\_\_} is the nearest thing, and it is not
  a value type & \ref{ch:objects} \\
\code{/} on two \code{int}s & \code{//}, which floors rather than
  truncating & \ref{ch:objects} \\
\code{decimal} & \api{decimal.Decimal} & \ref{ch:objects} \\
\code{Func<>}, a delegate & a function, annotated
  \code{Callable[[A], B]} & \ref{ch:functions} \\
an attribute a framework reads & nothing: a decorator \emph{runs} &
  \ref{ch:functions} \\
middleware round a call & a decorator & \ref{ch:functions} \\
\code{params}, optional, named & \code{*args}, defaults, keywords, and the
  \code{/} and \code{*} markers & \ref{ch:functions} \\
LINQ, deferred and re-runnable & a comprehension (eager) or a generator
  expression (deferred, single-pass) & \ref{ch:functions} \\
\code{yield return} & \code{yield}, plus \code{send} and a return value &
  \ref{ch:functions} \\
\code{using} / \code{IDisposable} & \code{with} /
  \api{contextlib.contextmanager} & \ref{ch:functions} \\
a \code{switch} expression & \code{match}, with sequence and mapping patterns &
  \ref{ch:functions} \\
\code{try} / \code{catch} / \code{finally} & \code{try} / \code{except} /
  \code{finally} / \code{else} & \ref{ch:errors} \\
\code{catch \{ \}} & \code{except Exception:} at most; a bare
  \code{except:} catches the interrupt too & \ref{ch:errors} \\
\code{throw;} & a bare \code{raise} & \ref{ch:errors} \\
\code{throw new X(..., inner)} & \code{raise X(...) from err} &
  \ref{ch:errors} \\
\code{AggregateException} & \code{ExceptionGroup} and \code{except*} &
  \ref{ch:errors} \\
\code{FooException} & \code{FooError} \dash{} ruff's N818 &
  \ref{ch:errors} \\
\code{TryGetValue} & ask forgiveness: \code{try} / \code{except KeyError}
  & \ref{ch:errors} \\
\code{using X;} & \code{import x} or \code{from x import name}, never
  \code{import *} & \ref{ch:imports} \\
an assembly & a package; a module is a file that runs once &
  \ref{ch:imports} \\
\code{internal} & a leading underscore, by convention only &
  \ref{ch:imports} \\
\code{IServiceCollection} & a composition root: a function,
  \api{functools.partial} and a \code{Protocol} & \ref{ch:imports} \\
\code{Program.cs} & the \code{[project.scripts]} entry point \code{uv
  init} writes & \ref{ch:imports} \\
\end{cheatsheet}

\section{Part III \dash{} Concurrency}
\label{sec:appA-part3}

\begin{cheatsheet}{C\#}{Python}{Ch.}
\code{Task} & a coroutine \dash{} but cold: calling it starts nothing &
  \ref{ch:asyncio} \\
\code{Task.Run} & \api{asyncio.to\_thread} & \ref{ch:asyncio} \\
\code{Task.WhenAll} & \api{asyncio.gather}, or \code{TaskGroup} if you
  want siblings cancelled & \ref{ch:asyncio} \\
\code{CancellationToken} & \code{CancelledError}, delivered at an
  \code{await}; nothing is polled & \ref{ch:asyncio} \\
\code{CancelAfter} & \api{asyncio.timeout} & \ref{ch:asyncio} \\
\code{ConfigureAwait(false)} & nothing: there is no synchronisation
  context & \ref{ch:asyncio} \\
\code{.Result}, \code{.Wait()} & nothing safe: a blocking call freezes the
  whole loop & \ref{ch:asyncio} \\
\code{SemaphoreSlim} & \api{asyncio.Semaphore} & \ref{ch:asyncio} \\
\code{IAsyncEnumerable<T>} & \code{async for} over an async generator &
  \ref{ch:asyncio} \\
\code{HttpClient} & \code{httpx.AsyncClient} \dash{} one instance, and it
  already has a deadline & \ref{ch:asyncio} \\
the thread pool grows & the default executor is capped, and the cap is
  small & \ref{ch:asyncio} \\
\end{cheatsheet}

\section{Part IV \dash{} Shipping}
\label{sec:appA-part4}

\begin{cheatsheet}{C\#}{Python}{Ch.}
ASP.NET Core & FastAPI on uvicorn & \ref{ch:services} \\
a controller & a router, and a function with a decorator &
  \ref{ch:services} \\
model binding & a pydantic model in the signature & \ref{ch:services} \\
\code{[ApiController]} 400 & a 422 carrying pydantic's error list &
  \ref{ch:services} \\
\code{IOptions<T>}, \code{appsettings.json} & \pkg{pydantic-settings} over
  the environment and \code{.env} & \ref{ch:services} \\
DI lifetimes & \code{Depends}, per-request; \code{yield} for a scope; the
  lifespan for singletons & \ref{ch:services} \\
middleware, first registered is outer & last registered is outer &
  \ref{ch:services} \\
Kestrel threads & uvicorn workers are processes, balanced per connection &
  \ref{ch:services} \\
\code{JsonPropertyName} & \code{serialization\_alias}; plain \code{alias}
  renames the input too & \ref{ch:services} \\
Entity Framework & SQLAlchemy & \ref{ch:data} \\
\code{DbContext} & \code{Session} \dash{} and it expires your objects on
  commit & \ref{ch:data} \\
\code{DbSet<T>}, LINQ & \code{select(T).where(...)}, which cannot take a
  lambda & \ref{ch:data} \\
\code{Include} & \code{selectinload}, against the same N+1 &
  \ref{ch:data} \\
\code{Add-Migration} & \code{alembic revision --autogenerate}, which you
  must read & \ref{ch:data} \\
\code{Update-Database} & \code{alembic upgrade head} & \ref{ch:data} \\
a DataTable & a \pkg{polars} frame & \ref{ch:data} \\
xUnit & pytest & \ref{ch:testing} \\
\code{[Fact]}, \code{[Theory]} & a \code{test\_} function,
  \code{@pytest.mark.parametrize} & \ref{ch:testing} \\
\code{IClassFixture<T>} & a fixture function with a scope &
  \ref{ch:testing} \\
Moq & \api{unittest.mock}, patched where the name is looked up &
  \ref{ch:testing} \\
\code{Assert.Equal} & a plain \code{assert}, which pytest rewrites to
  explain itself & \ref{ch:testing} \\
FluentAssertions & the plain \code{assert}, again & \ref{ch:testing} \\
\code{ILogger<T>} & \pkg{structlog} over the standard \api{logging}
  module & \ref{ch:observability} \\
\code{AsyncLocal<T>} & \api{contextvars} & \ref{ch:observability} \\
OpenTelemetry & OpenTelemetry, and the same spans & \ref{ch:observability} \\
\code{IHostedService} shutdown & uvicorn's lifespan, which waits for ever
  by default & \ref{ch:observability} \\
\code{dotnet publish}, a runtime image & a multi-stage Dockerfile and
  \code{uv sync --frozen --no-dev} & \ref{ch:observability} \\
\code{dotnet list package -{}-vulnerable} & \code{pip-audit}, over a
  \code{pylock.toml} you export & \ref{ch:observability} \\
\end{cheatsheet}

\section{Part V \dash{} Python for AI work}
\label{sec:appA-part5}

\begin{cheatsheet}{C\#}{Python}{Ch.}
an SDK client & a pydantic model, an httpx client and a streaming
  iterator & \ref{ch:aitoolkit} \\
\code{JsonSerializer.Deserialize} & \code{model\_validate\_json}, and it
  is faster than parsing first & \ref{ch:aitoolkit} \\
a strict deserialiser & \code{strict=True}, which costs nothing &
  \ref{ch:aitoolkit} \\
\code{IAsyncEnumerable} over a stream & \code{async for} over the SDK's
  stream & \ref{ch:aitoolkit} \\
a scratch console project & a notebook, and the book says when one is
  right & \ref{ch:aitoolkit} \\
xUnit over a recorded trace & \code{pytest} over a \code{Trace}, with
  twelve deterministic assertions & \ref{ch:traceassert} \\
\code{dotnet pack}, \code{nuget push} & \code{uv build}, \code{uv
  publish} with trusted publishing & \ref{ch:traceassert} \\
\end{cheatsheet}
